\documentclass[../main.tex]{subfiles}
\begin{document}

\section{Problema Lungcircuit (Lot 2025)}
\label{problem:lungcircuit}

\href{https://kilonova.ro/problems/3807}{enunț}
$\bullet$
\hyperref[code:lungcircuit]{surse}

Problema Lungcircuit îmi aparține. Inițial căutam o problemă pentru lotul de juniori, bazată pe jocul \href{https://ruwix.com/twisty-puzzles/hungarian-rings-sliding-puzzle/}{Hungarian Rings}. A rezultat o problemă prea grea pentru juniori. În seara dinaintea unei probe, comisia de seniori mi-a cerut problema ca să o dea a doua zi, pentru că problema avea un mare avantaj: era la cheie (teste, subtaskuri, soluții pentru toate subtaskurile, chiar și editorial). Uneori gradul de pregătire bate calitatea problemei. \emoji{smiling-face-with-tear}

Problema are o soluție „tractor” și o soluție elegantă, care este mai mult de idee. \href{https://drive.google.com/drive/folders/1EqiZNtAgVsJ-FbrPRqzF-OIpyOPeztnR}{Editorialul} include și soluții pentru subtaskuri.

\subsection{Soluția „tractor”}

Introducem doi termeni pentru a ușura descrierea. Numim pozițiile comune \textbf{noduri}. Primul nod este la poziția 0, iar ultimul nod este la poziția $(k-1)d$. Numim \textbf{lanțuri exterioare} pozițiile $(k-1)d + 1 \dots n - 1$ pe fiecare circuit.

Pentru subtaskurile mici, putem menține circuitele complete la fiecare mișcare. Pentru o soluție în $\bigoh((q + n) \log k)$, să schimbăm paradigma. Să urmărim traseul fiecărei bile și să aflăm poziția ei finală după cele $q$ mutări. Vom accelera simularea procesînd rapid porțiuni de circuit. Pornind de la mutarea 0, fiecare bilă va trece prin următoarele etape (sau o submulțime a lor):

\begin{enumerate}
  \item Avansează pînă la următorul nod.
  \item Avansează pînă la ultimul nod.
  \item Avansează pînă la primul nod (parcurgînd lanțul exterior).
  \item Găsește ultima mutare cînd bila mai revine în primul nod.
  \item Avansează un număr întreg de noduri (cel mult $k-1$).
  \item Avansează pozițiile rămase pînă la epuizarea operațiilor.
\end{enumerate}

Pentru (1), (3) și (6), știind că bila se află pe circuitul $c$ la mutarea $t$ și că dorim să avanseze $p$ pași, trebuie să putem răspunde la întrebarea „ce indice are a $p$-a mutare de tip $c$, începînd cu mutarea $t$”? Putem răspunde la aceste întrebări cu o precalculare ușoară, colectînd în doi vectori pozițiile literelor \textbf{R} și \textbf{A} în șirul de mutări.

Pentru (2) și (5), să remarcăm că bila pornește dintr-un nod, deci prima mișcare dictează pe ce drum pornește bila. Pentru a avansa un interval în $\bigoh(1)$, dacă prima mișcare este $c$ la momentul $t$, trebuie să găsim rapid a $d$-a apariție a lui $c$, începînd cu mutarea $t$, în șirul de mutări. Și această informație poate fi precalculată.

Noi dorim să avansăm eficient $\bigoh(k)$ intervale, deci peste informația de mai sus construim o tabelă rară (\textit{sparse table}). Astfel putem avansa pînă la ultimul nod în $\bigoh(\log k)$.

Pentru (4), dorim să răspundem la întrebări de forma: „dacă o bilă se află în nodul 0 la momentul $t$, cînd va reveni ea în nodul 0?”. Putem refolosi exact codul anterior. Dacă răspunsul există și este $t'$, atunci putem completa de la sfîrșitul spre începutul vectorului de mutări momentul \textbf{ultimei} reveniri în nodul 0: el va fi același pentru $t$ și pentru $t'$.

\subsection{Soluție cu liste}

O soluție în $\bigoh(n + q)$, semnificativ mai simplu de codat, pornește de la următoarea observație. Bilele din noduri (în exemplul dat, 0, 3, 6, 9 și 12) au o traiectorie similară. Dacă prima mutare este \textbf{R} și bila 0 intră pe circuitul roșu, la fel vor face și celelalte. Cînd 0 revine într-un nod, la fel și celelalte. Similar sînt unite și bilele 1, 4, 7, 10 și 13.

De aceea, putem să grupăm aceste bile și să le mutăm pe toate cu efort mic. Concret, reprezentăm un circuit ca:

\begin{itemize}
  \item $d$ liste, prima de $k$ bile, iar restul de cîte $k-1$ bile. Numim această structură \textbf{rețea}, căci elementele listelor sînt intercalate. Prima listă conține nodurile și este comună celor două circuite.
  \item Lanțul exterior (inițial 16-31, respectiv 42-57).
\end{itemize}

Ce se întîmplă la o mutare \textbf{R} (cazul \textbf{A} este simetric)?

\begin{itemize}
  \item Lanțul exterior capătă o bilă la început, și anume ultimul nod (15).
  \item Lanțul exterior pierde o bilă de la final (31), care ajunge la începutul ultimei liste din rețea, care devine 31, 2, 5, 8, 11, 14.
  \item Permutăm circular listele rețelei. Ultima listă devine prima. Astfel, bila 31, care acum stă în nodul 0, devine în mod corect prima bilă din prima listă.
  \item Notificăm circuitul albastru că prima sa listă nu mai este $0, 3, \dots$, ci este $31, 2, \dots$
\end{itemize}

Cu această reprezentare putem simula toți pașii de mai sus, și deci toate mutările, în $\bigoh(1)$ fiecare. Putem reprezenta atît listele din rețea, cît și lista de liste (rețeaua însăși), ca pe buffere circulare, ca să le putem permuta circular în $\bigoh(1)$.

Unele implementări se simplifică considerabil dacă stocăm ultimul nod separat de rețea. Atunci toate listele din rețea au lungime egală, $k-1$.

Implementarea cu \ccode{deque}-uri este chiar mai scurtă, dar este foarte ineficientă cînd $d$ este mare, deoarece creează o mulțime de \ccode{deque}-uri mici.

\end{document}
